\begin{center}
\thispagestyle{empty}
\includegraphics[height=4cm]{HojaTitulo/EscudoUN.pdf}\\[2.6cm]
\textbf{\huge
Topologías de Alexandroff sobre grafos simples}\\[4.2cm]
\Large\textbf{Juan Camilo Lozano Suárez}\\[4.2cm]
\small Universidad Nacional de Colombia\\
Facultad de Ciencias\\
Departamento de Matemáticas\\
Bogotá, Colombia\\
2026\\
\end{center}

\newpage{\pagestyle{empty}\cleardoublepage}

\newpage
\begin{center}
\thispagestyle{empty} \vspace*{0cm} \textbf{\huge
Topologías de Alexandroff sobre grafos simples}\\[2.5cm]
\Large\textbf{Juan Camilo Lozano Suárez}\\[2.5cm]
\small Trabajo de grado presentado como requisito parcial para optar al título de:\\
\textbf{Matemático}\\[2cm]
Directora:\\
\textbf{Ibeth Marcela Rubio Perilla}\\[1.8cm]
Líneas de Investigación:\\
Topología\\
Teoría de grafos\\[2cm]
Universidad Nacional de Colombia\\
Facultad de Ciencias\\
Departamento de Matemáticas\\
Bogotá, Colombia\\
2026\\
\end{center}
% ---------------- Página en blanco ----------------
\clearpage
\thispagestyle{empty}
\null
\clearpage

% ---------------- Dedicatoria ----------------
\thispagestyle{empty}
\begin{flushright}
   \begin{minipage}{0.6\textwidth}
      \raggedleft
      \textit{A mis padres, por su apoyo, paciencia y comprensión, que no tienen cardinal.}
   \end{minipage}
\end{flushright}
\clearpage

% ---------------- Página en blanco ----------------
\thispagestyle{empty}
\null
\clearpage% ---------------- Epígrafes ----------------
\clearpage
\thispagestyle{empty}
\begin{flushright}
   \begin{minipage}{0.6\textwidth}
      \raggedleft
      \textit{Para sobrevivir al horror hay que inventar algún tipo de belleza.}\par
      \vspace{1ex}
      ---Fernando Zalamea\footnotemark
   \end{minipage}
\end{flushright}
\footnotetext{Refiriéndose a la invención del concepto de \textit{haz} por parte de Jean Leray, en 1942, cuando este se encontraba recluido en un campo de concentración \cite{zalamea2018}.}

\vspace{3cm}

\begin{flushright}
   \begin{minipage}{0.6\textwidth}
      \raggedleft
      \textit{Nadie puede leer dos mil libros. En los cuatro siglos que vivo no habré pasado de media docena.}\par
      \vspace{1ex}
      ---Jorge Luis Borges, \textit{Utopía de un hombre que está cansado}
   \end{minipage}
\end{flushright}
\vspace*{\fill}
\clearpage
% ---------------- Página en blanco ----------------
\clearpage
\thispagestyle{empty}
\null
\clearpage

\newpage
\textbf{\LARGE Resumen}\\\\
El resumen es una presentaci\'{o}n abreviada y precisa (la NTC 1486 de 2008 recomienda revisar la norma ISO 214 de 1976). Se debe usar una extensi\'{o}n m\'{a}xima de 12 renglones. Se recomienda que este resumen sea anal\'{\i}tico, es decir, que sea completo, con informaci\'{o}n cuantitativa y cualitativa, generalmente incluyendo los siguientes aspectos: objetivos, dise\~{n}o, lugar y circunstancias, pacientes (u objetivo del estudio), intervenci\'{o}n, mediciones y principales resultados, y conclusiones. Al final del resumen se deben usar palabras claves tomadas del texto (m\'{\i}nimo 3 y m\'{a}ximo 7 palabras), las cuales permiten la recuperaci\'{o}n de la informaci\'{o}n.\\

\textbf{\small Palabras clave: (m\'{a}ximo 10 palabras, preferiblemente seleccionadas de las listas internacionales que permitan el indizado cruzado)}.\\

\textbf{\LARGE Abstract}\\\\
Es el mismo resumen pero traducido al ingl\'{e}s. Se debe usar una extensi\'{o}n m\'{a}xima de 12 renglones. Al final del Abstract se deben traducir las anteriores palabras claves tomadas del texto (m\'{\i}nimo 3 y m\'{a}ximo 7 palabras), llamadas keywords. Es posible incluir el resumen en otro idioma diferente al espa\~{n}ol o al ingl\'{e}s, si se considera como importante dentro del tema tratado en la investigaci\'{o}n, por ejemplo: un trabajo dedicado a problemas ling\"{u}\'{\i}sticos del mandar\'{\i}n seguramente estar\'{\i}a mejor con un resumen en mandar\'{\i}n.\\

\textbf{\small Keywords: palabras clave en ingl\'{e}s(m\'{a}ximo 10 palabras, preferiblemente seleccionadas de las listas internacionales que permitan el indizado cruzado)}\\
